\clearpage
\phantomsection\addcontentsline{toc}{chapter}{References}
\renewcommand{\bibname}{References}
\begingroup\singlespacing\footnotesize
\begin{thebibliography}{99}
\bibitem{r1} A. Radford, J. Wu, R. Child, D. Luan, D. Amodei, and I. Sutskever, ``Language models are unsupervised multitask learners,'' \textit{OpenAI Technical Report}, 2019.

\bibitem{r2} K. Meng, D. Bau, A. Andonian, and Y. Belinkov, ``Locating and editing factual associations in GPT,'' in \textit{Advances in Neural Information Processing Systems (NeurIPS)}, vol. 35, 2022, pp. 17359--17372.

\bibitem{r3} K. Meng, A. S. Sharma, A. Andonian, Y. Belinkov, and D. Bau, ``Mass-editing memory in a transformer,'' in \textit{Proc. International Conference on Learning Representations (ICLR)}, 2023.

\bibitem{r4} A. Gupta, A. Rao, and G. Anumanchipalli, ``Model editing at scale leads to gradual and catastrophic forgetting,'' in \textit{Findings of the Association for Computational Linguistics: ACL 2024}, 2024; arXiv:2401.07453.

\bibitem{r5} N. Elhage, T. Hume, C. Olsson, \textit{et al.}, ``Toy models of superposition,'' \textit{Transformer Circuits Thread}, 2022.

\bibitem{r6} T. Bricken, A. Templeton, J. Batson, \textit{et al.}, ``Towards monosemanticity: Decomposing language models with dictionary learning,'' \textit{Transformer Circuits Thread}, 2023.

\bibitem{r7} H. Cunningham, A. Ewart, L. Riggs, R. Huben, and L. Sharkey, ``Sparse autoencoders find highly interpretable features in language models,'' in \textit{Proc. International Conference on Learning Representations (ICLR)}, 2024.

\bibitem{r8} A. Templeton, T. Conerly, \textit{et al.}, ``Scaling monosemanticity: Extracting interpretable features from Claude 3 Sonnet,'' \textit{Transformer Circuits Thread}, 2024.

\bibitem{r9} A. M. Turner, L. Thiergart, G. Leech, D. Udell, J. J. Vazquez, U. Mini, and M. MacDiarmid, ``Steering language models with activation engineering,'' \textit{arXiv preprint arXiv:2308.10248}, 2023.

\bibitem{r10} N. Rimsky, N. Gabrieli, J. Schulz, M. Tong, E. Hubinger, and A. M. Turner, ``Steering Llama 2 via contrastive activation addition,'' in \textit{Proc. 62nd Annual Meeting of the Association for Computational Linguistics (ACL)}, 2024, pp. 15504--15522.

\bibitem{r11} A. Zou, L. Phan, S. Chen, \textit{et al.}, ``Representation engineering: A top-down approach to AI transparency,'' \textit{arXiv preprint arXiv:2310.01405}, 2023.

\bibitem{r12} Z. Wu, A. Arora, A. Geiger, Z. Wang, J. Huang, D. Jurafsky, C. D. Manning, and C. Potts, ``AxBench: Steering LLMs? Even simple baselines outperform sparse autoencoders,'' in \textit{Proc. International Conference on Machine Learning (ICML)}, 2025; arXiv:2501.17148.

\bibitem{r13} N. Nanda and J. Bloom, ``TransformerLens,'' GitHub repository, 2022. [Online]. Available: \url{https://github.com/TransformerLensOrg/TransformerLens}

\bibitem{r14} J. Bloom, C. Tigges, A. Duong, and D. Chanin, ``SAELens,'' GitHub repository, 2024. [Online]. Available: \url{https://github.com/decoderesearch/SAELens}

\bibitem{r15} J. Bloom, ``Open source sparse autoencoders for all residual stream layers of GPT2 small,'' \textit{AI Alignment Forum}, 2024.

\bibitem{r16} J. Kram\'ar, T. Lieberum, R. Shah, and N. Nanda, ``AtP*: An efficient and scalable method for localizing LLM behaviour to components,'' \textit{arXiv preprint arXiv:2403.00745}, 2024.

\bibitem{r19} T. Lieberum, S. Rajamanoharan, A. Conmy, \textit{et al.}, ``Gemma Scope: Open sparse autoencoders everywhere all at once on Gemma 2,'' \textit{arXiv preprint arXiv:2408.05147}, 2024.

\bibitem{r20} J. Lin and J. Bloom, ``Neuronpedia: Interactive reference and tooling for analyzing neural networks,'' 2023. [Online]. Available: \url{https://www.neuronpedia.org}

\bibitem{r21} S. Heimersheim and N. Nanda, ``How to use and interpret activation patching,'' \textit{arXiv preprint arXiv:2404.15255}, 2024.

\bibitem{r22} A. Paszke, S. Gross, F. Massa, \textit{et al.}, ``PyTorch: An imperative style, high-performance deep learning library,'' in \textit{Advances in Neural Information Processing Systems (NeurIPS)}, vol. 32, 2019, pp. 8024--8035.

\bibitem{r23} D. Arad, A. Mueller, and Y. Belinkov, ``SAEs are good for steering -- if you select the right features,'' in \textit{Proc. Conference on Empirical Methods in Natural Language Processing (EMNLP)}, 2025; arXiv:2505.20063.

\bibitem{r24} L. Smith, S. Rajamanoharan, A. Conmy, C. McDougall, T. Lieberum, J. Kram\'ar, R. Shah, and N. Nanda, ``Negative results for SAEs on downstream tasks and deprioritising SAE research (GDM mech interp team progress update \#2),'' \textit{LessWrong}, Mar. 2025. [Online]. Available: \url{https://www.lesswrong.com/posts/4uXCAJNuPKtKBsi28/negative-results-for-saes-on-downstream-tasks}

\bibitem{r25} M. Chaudhary and A. Geiger, ``Evaluating open-source sparse autoencoders on disentangling factual knowledge in GPT-2 small,'' \textit{arXiv preprint arXiv:2409.04478}, 2024.

\bibitem{r26} M. Cui, L. Shen, and X. Yang, ``SAE interventions are unreliable: Post-intervention recovery of suppressed behavior,'' \textit{arXiv preprint arXiv:2606.18322}, 2026.

\bibitem{r27} E. Farrell, Y.-T. Lau, and A. Conmy, ``Applying sparse autoencoders to unlearn knowledge in language models,'' \textit{arXiv preprint arXiv:2410.19278}, 2024.

\bibitem{r28} S. Chalnev, M. Siu, and A. Conmy, ``Improving steering vectors by targeting sparse autoencoder features,'' \textit{arXiv preprint arXiv:2411.02193}, 2024.

\bibitem{r29} Y. Zhao, A. Devoto, G. Hong, X. Du, A. P. Gema, H. Wang, X. He, K.-F. Wong, and P. Minervini, ``Steering knowledge selection behaviours in LLMs via SAE-based representation engineering,'' in \textit{Proc. NAACL}, 2025; arXiv:2410.15999.

\bibitem{r30} P. Hase, M. Bansal, B. Kim, and A. Ghandeharioun, ``Does localization inform editing? Surprising differences in causality-based localization vs.\ knowledge editing in language models,'' in \textit{Advances in Neural Information Processing Systems (NeurIPS)}, vol. 36, 2023.

\bibitem{r31} C. Zheng, L. Li, Q. Dong, Y. Fan, Z. Wu, J. Xu, and B. Chang, ``Can we edit factual knowledge by in-context learning?'' in \textit{Proc. Conference on Empirical Methods in Natural Language Processing (EMNLP)}, 2023; arXiv:2305.12740.

\bibitem{r32} P. Wang, N. Zhang, B. Tian, \textit{et al.}, ``EasyEdit: An easy-to-use knowledge editing framework for large language models,'' in \textit{Proc. 62nd Annual Meeting of the Association for Computational Linguistics (System Demonstrations)}, 2024; arXiv:2308.07269.

\bibitem{r33} S. Marks, C. Rager, E. J. Michaud, Y. Belinkov, D. Bau, and A. Mueller, ``Sparse feature circuits: Discovering and editing interpretable causal graphs in language models,'' in \textit{Proc. International Conference on Learning Representations (ICLR)}, 2025; arXiv:2403.19647.

\bibitem{r34} P. Lewis, E. Perez, A. Piktus, \textit{et al.}, ``Retrieval-augmented generation for knowledge-intensive NLP tasks,'' in \textit{Advances in Neural Information Processing Systems (NeurIPS)}, vol. 33, 2020.

\bibitem{r36} D. P. Kingma and J. Ba, ``Adam: A method for stochastic optimization,'' in \textit{Proc. International Conference on Learning Representations (ICLR)}, 2015.

\bibitem{r37} nostalgebraist, ``Interpreting GPT: the logit lens,'' \textit{LessWrong}, Aug. 2020. [Online]. Available: \url{https://www.lesswrong.com/posts/AcKRB8wDpdaN6v6ru/interpreting-gpt-the-logit-lens}

\bibitem{r39} Q. McNemar, ``Note on the sampling error of the difference between correlated proportions or percentages,'' \textit{Psychometrika}, vol. 12, no. 2, pp. 153--157, 1947.

\bibitem{r40} F. Wilcoxon, ``Individual comparisons by ranking methods,'' \textit{Biometrics Bulletin}, vol. 1, no. 6, pp. 80--83, 1945.

\bibitem{r41} K. Lenc and A. Vedaldi, ``Understanding image representations by measuring their equivariance and equivalence,'' in \textit{Proc. IEEE Conference on Computer Vision and Pattern Recognition (CVPR)}, 2015.

\bibitem{r42} Y. Bansal, P. Nakkiran, and B. Barak, ``Revisiting model stitching to compare neural representations,'' in \textit{Advances in Neural Information Processing Systems (NeurIPS)}, vol. 34, 2021.

\bibitem{r43} S. Kornblith, M. Norouzi, H. Lee, and G. Hinton, ``Similarity of neural network representations revisited,'' in \textit{Proc. International Conference on Machine Learning (ICML)}, 2019.

\bibitem{r44} A. Chen, J. Merullo, A. Stolfo, and E. Pavlick, ``Transferring linear features across language models with model stitching,'' in \textit{Advances in Neural Information Processing Systems (NeurIPS)}, 2025; arXiv:2506.06609.

\bibitem{r45} A. E. Hoerl and R. W. Kennard, ``Ridge regression: Biased estimation for nonorthogonal problems,'' \textit{Technometrics}, vol. 12, no. 1, pp. 55--67, 1970.

\bibitem{r46} B. Lester, R. Al-Rfou, and N. Constant, ``The power of scale for parameter-efficient prompt tuning,'' in \textit{Proc. Conference on Empirical Methods in Natural Language Processing (EMNLP)}, 2021.

\bibitem{r47} X. L. Li and P. Liang, ``Prefix-tuning: Optimizing continuous prompts for generation,'' in \textit{Proc. 59th Annual Meeting of the Association for Computational Linguistics (ACL)}, 2021.

\bibitem{r48} N. Subramani, N. Suresh, and M. E. Peters, ``Extracting latent steering vectors from pretrained language models,'' in \textit{Findings of the Association for Computational Linguistics: ACL 2022}, 2022, pp. 566--581.

\bibitem{r49} J. Dunefsky and A. Cohan, ``One-shot optimized steering vectors mediate safety-relevant behaviors in LLMs,'' in \textit{Conference on Language Modeling (COLM)}, 2025; arXiv:2502.18862.

\bibitem{r50} C. Szegedy, W. Zaremba, I. Sutskever, J. Bruna, D. Erhan, I. Goodfellow, and R. Fergus, ``Intriguing properties of neural networks,'' in \textit{Proc. International Conference on Learning Representations (ICLR)}, 2014.

\bibitem{r51} I. J. Goodfellow, J. Shlens, and C. Szegedy, ``Explaining and harnessing adversarial examples,'' in \textit{Proc. International Conference on Learning Representations (ICLR)}, 2015.

\bibitem{r52} C. Zhang, S. Bengio, M. Hardt, B. Recht, and O. Vinyals, ``Understanding deep learning requires rethinking generalization,'' in \textit{Proc. International Conference on Learning Representations (ICLR)}, 2017.

\bibitem{r53} N. Elhage, N. Nanda, C. Olsson, \textit{et al.}, ``A mathematical framework for transformer circuits,'' \textit{Transformer Circuits Thread}, 2021.

\bibitem{r54} S. Kullback and R. A. Leibler, ``On information and sufficiency,'' \textit{The Annals of Mathematical Statistics}, vol. 22, no. 1, pp. 79--86, 1951.

\bibitem{r55} S. Merity, C. Xiong, J. Bradbury, and R. Socher, ``Pointer sentinel mixture models,'' in \textit{Proc. International Conference on Learning Representations (ICLR)}, 2017.

\bibitem{r56} W. J. Dixon and A. M. Mood, ``The statistical sign test,'' \textit{Journal of the American Statistical Association}, vol. 41, no. 236, pp. 557--566, 1946.

\bibitem{r57} B. Efron, ``Bootstrap methods: Another look at the jackknife,'' \textit{The Annals of Statistics}, vol. 7, no. 1, pp. 1--26, 1979.

\bibitem{r58} J. Vig, S. Gehrmann, Y. Belinkov, S. Qian, D. Nevo, Y. Singer, and S. Shieber, ``Investigating gender bias in language models using causal mediation analysis,'' in \textit{Advances in Neural Information Processing Systems (NeurIPS)}, vol. 33, 2020.

\bibitem{r59} J. Pearl, ``Direct and indirect effects,'' in \textit{Proc. 17th Conference on Uncertainty in Artificial Intelligence (UAI)}, 2001, pp. 411--420.

\bibitem{r60} M. Geva, R. Schuster, J. Berant, and O. Levy, ``Transformer feed-forward layers are key-value memories,'' in \textit{Proc. Conference on Empirical Methods in Natural Language Processing (EMNLP)}, 2021, pp. 5484--5495.

\bibitem{r61} A. Holtzman, J. Buys, L. Du, M. Forbes, and Y. Choi, ``The curious case of neural text degeneration,'' in \textit{Proc. International Conference on Learning Representations (ICLR)}, 2020.

\bibitem{r62} J. Li, M. Galley, C. Brockett, J. Gao, and B. Dolan, ``A diversity-promoting objective function for neural conversation models,'' in \textit{Proc. Conference of the North American Chapter of the Association for Computational Linguistics (NAACL)}, 2016, pp. 110--119.

\bibitem{r63} T. Wolf, L. Debut, V. Sanh, \textit{et al.}, ``Transformers: State-of-the-art natural language processing,'' in \textit{Proc. EMNLP: System Demonstrations}, 2020, pp. 38--45.

\bibitem{r64} K. Meng \textit{et al.}, ``ROME: Rank-One Model Editing, official implementation,'' GitHub repository, commit 0874014 (MIT licence). [Online]. Available: \url{https://github.com/kmeng01/rome}

\bibitem{r65} K. Meng, D. Bau, A. Andonian, and Y. Belinkov, ``Locating and editing factual associations in GPT: project page, CounterFact and known-facts data.'' [Online]. Available: \url{https://rome.baulab.info}

\bibitem{r66} Streamlit Inc., ``Streamlit: A faster way to build and share data apps,'' software. [Online]. Available: \url{https://streamlit.io}

\medskip\noindent\textbf{Project records.} Results that come from our own experiments are cited to the records below (labels P1 to P19). All are files in the project repository (branch \texttt{cross-model-transfer}); every figure and table in this report is drawn from them (Appendix~D).\par\smallskip

\bibitem[P1]{p1} FeatureScalpel project, ``Research Journal, entry 1: project progress and roadmap (first prototype, causal selector, eligibility rule),'' July 2026. \texttt{docs/Research\_Journal/1.md}.

\bibitem[P2]{p2} FeatureScalpel project, ``Research Journal, entries 2 and 3: iterative ablation of the target's helpers and of the competitor's drivers,'' 19--20 July 2026. \texttt{docs/Research\_Journal/2.md, 3.md}.

\bibitem[P3]{p3} FeatureScalpel project, ``Research Journal, entries 4 to 8: whole-combination safety, weighted multi-competitor muting, mute-and-boost diagnostics and parameter sweep,'' 24--25 July 2026. \texttt{docs/Research\_Journal/4.md} to \texttt{8.md}.

\bibitem[P4]{p4} FeatureScalpel project, ``Research Journal, entries 9 to 11: SAE dictionary audit, layer-depth benchmark, benchmark infrastructure and stage profiling,'' 4--6 August 2026. \texttt{docs/Research\_Journal/9.md} to \texttt{11.md}; data \texttt{benchmark\_results/layer\_benchmark\_*.json}, \texttt{benchmark\_results/layer\_benchmark\_20260909\_summary.csv}.

\bibitem[P5]{p5} FeatureScalpel project, ``Research Journal, entries 12 to 15: delta patching, hardware comparison, execution-cost audit, GPU batching of candidate evaluation,'' 9 August to 18 September 2026. \texttt{docs/Research\_Journal/12.md} to \texttt{15.md}.

\bibitem[P6]{p6} FeatureScalpel project, ``Research Journal, entries 16 to 19: best-so-far tracking, cumulative sweep, combined hook, pool refill,'' 18--24 September 2026. \texttt{docs/Research\_Journal/16.md} to \texttt{19.md}.

\bibitem[P7]{p7} FeatureScalpel project, ``Research Journal, entry 20: candidate source, last token against all prompt positions,'' 25--26 September 2026. \texttt{docs/Research\_Journal/20.md}; data \texttt{benchmark\_results/candidate\_source\_study/}.

\bibitem[P8]{p8} FeatureScalpel project, ``Safety-filter study (journal entry 21 placeholder and result packs),'' 26 September 2026. \texttt{docs/Research\_Journal/packs/entry21\_*}; specifications \texttt{data/safety\_filter\_spec\_*.json}.

\bibitem[P9]{p9} FeatureScalpel project, ``Research Journal, entry 22: do later layers undo the edit, is the SAE the limit,'' 30 September 2026. \texttt{docs/Research\_Journal/22.md}; data \texttt{docs/Research\_Journal/packs/repair\_diagnostics/}.

\bibitem[P10]{p10} FeatureScalpel project, ``Research Journal, entries 23 and 24: CounterFact hard set, split, per-prompt cache and headroom measurement,'' 1--2 October 2026. \texttt{docs/Research\_Journal/23.md, 24.md}; data \texttt{data/counterfact\_*.json}, \texttt{docs/Research\_Journal/packs/headroom/}.

\bibitem[P11]{p11} FeatureScalpel project, ``Research Journal, entry 25: literature check and assessment of the research direction,'' 1 October 2026. \texttt{docs/Research\_Journal/25.md}.

\bibitem[P12]{p12} FeatureScalpel project, ``Research Journal, entries 26 and 27: learned strengths (PromptNet, FeatureNet),'' 2--3 October 2026. \texttt{docs/Research\_Journal/26.md, 27.md}; data \texttt{docs/Research\_Journal/packs/strength\_models/}.

\bibitem[P13]{p13} FeatureScalpel project, ``Research Journal, entries 28 and 29: gradient-descent editing and the equal-budget comparison with the sweep,'' 3 October 2026. \texttt{docs/Research\_Journal/28.md, 29.md}; data \texttt{docs/Research\_Journal/packs/gradient\_descent/}, \texttt{packs/sweep\_vs\_gradient/}, \texttt{images/e29\_figure\_data.json}.

\bibitem[P14]{p14} FeatureScalpel project, ``Research Journal, entry 30: additive edits, random-target control, reworded-prompt, generation and instruction tests, Prototype lab,'' 4 October 2026. \texttt{docs/Research\_Journal/30.md}; data \texttt{docs/Research\_Journal/packs/additive\_control/}, \texttt{packs/generalisation/}, \texttt{packs/generation\_quality/}.

\bibitem[P15]{p15} FeatureScalpel project (branch \texttt{rome-factual-editing}), ``ROME baseline: progress log, causal-tracing summary and edit-layer sweep on GPT-2 Small,'' 7 October 2026. \texttt{docs/rome\_baseline/PROGRESS.md}, \texttt{data/comparison/trace\_summary.json}, \texttt{data/comparison/gpt2\_knowns.json}, \texttt{benchmark\_results/rome\_vs\_scalpel/layer\_sweep\_dev.csv}; code \texttt{third\_party/rome/}, \texttt{src/rome\_baseline/}, \texttt{scripts/rome/}.

\bibitem[P16]{p16} FeatureScalpel project, ``Research Journal, entries 31 to 33: translator between GPT-2 Small and Medium, country swap, controlled Export, search for a better translator,'' 9--10 October 2026. \texttt{docs/Research\_Journal/31.md} to \texttt{33.md}; data \texttt{docs/Research\_Journal/packs/cross\_model\_transfer/}.

\bibitem[P17]{p17} FeatureScalpel project, ``Research Journal, entries 34 to 37: Import, varied-edit training, softer Import, SAE-basis control,'' 10--11 October 2026. \texttt{docs/Research\_Journal/34.md} to \texttt{37.md}; data \texttt{docs/Research\_Journal/packs/cross\_model\_transfer/}.

\bibitem[P18]{p18} FeatureScalpel project, ``Cross-model transfer plan with pass bars (gates) written before each run,'' October 2026. \texttt{docs/cross\_model\_transfer/PLAN.md}.

\bibitem[P19]{p19} FeatureScalpel project, ``Research roadmap, handoff notes and architecture notes.'' \texttt{docs/research\_roadmap.md}, \texttt{docs/handoff\_*.md}, \texttt{docs/architecture.md}.
\end{thebibliography}
\endgroup
